\documentclass[document/type=manual,localization/language=english]{insr}

\INSRConfigure{
  metadata/title = {INSR~LaTeX~Package~Manual},
  metadata/subtitle = {Public~v4~class~and~runtime},
  metadata/short-title = {INSR~Manual},
  metadata/author = {INSR~Research~Team},
  metadata/institution = {INSR~Research~Consortium}
}

\begin{document}

\INSRMakeTitle
\tableofcontents

\INSRChapter{Package Overview}

The INSR LaTeX package provides a LuaLaTeX-first document system for interdisciplinary trauma research. It includes scientific-paper, presentation, and clinical-manual adapters, Overleaf-compatible build settings, and GitHub CI integration.

\INSRChapter{Configuration}

Document type, design, localization, metadata, and bibliography resources are configured through the public INSR key-value API. Productive documents should use the canonical root class rather than deprecated legacy classes.

\begin{itemize}
  \item Select the output through \texttt{document/type}.
  \item Select language through \texttt{localization/language}.
  \item Keep project-wide metadata in \texttt{config/project-config.tex}.
  \item Keep document content outside the class implementation.
\end{itemize}

\INSRChapter{Clinical Review Helpers}

\begin{itemize}
  \item Document clinical review items explicitly.
  \item Document statistical-review items separately.
  \item Document technical-review items separately.
  \item Preserve human oversight for every safety-relevant decision.
\end{itemize}

\end{document}
